% TODO make hw stylesheet
\documentclass[10pt]{article}
\usepackage[letterpaper,top=1in,left=1in,right=1in]{geometry}
\usepackage[dvipsnames,svgnames]{xcolor}
\usepackage[shortlabels]{enumitem}
\usepackage[framemethod=TikZ]{mdframed}
\usepackage{amsmath,amssymb,amsthm}
\usepackage[colorlinks]{hyperref}
\usepackage{multicol}
\usepackage{tabularx}
\usepackage{bbm}

\hypersetup{
    colorlinks=true,
    linkcolor=blue,
    filecolor=magenta,
    urlcolor=magenta,
    pdftitle={MAT 549 HW}
}

\usepackage{mathtools}
\usepackage{thmtools}
\usepackage[textsize=footnotesize]{todonotes}
\usepackage{titling}
\usepackage{cancel}

\reversemarginpar
\mdfdefinestyle{mdbluebox}{roundcorner=10pt,innerbottommargin=9pt,
linecolor=blue,backgroundcolor=TealBlue!5,}
\declaretheoremstyle[headfont=\sffamily\bfseries\color{MidnightBlue},
mdframed={style=mdbluebox},]{thmbluebox}

\mdfdefinestyle{mdbloobox}{frametitlefont=\bfseries,innerbottommargin=8pt,
nobreak=true,backgroundcolor=TealBlue!5,linecolor=blue,}
\declaretheoremstyle[headfont=\bfseries\color{MidnightBlue},
mdframed={style=mdbloobox},headpunct={\\[3pt]},postheadspace=0pt,]{thmbloobox}

\mdfdefinestyle{mdredbox}{frametitlefont=\bfseries,innerbottommargin=8pt,
nobreak=true,backgroundcolor=Salmon!5,linecolor=RawSienna,}
\declaretheoremstyle[headfont=\bfseries\color{RawSienna},
mdframed={style=mdredbox},headpunct={\\[3pt]},postheadspace=0pt,]{thmredbox}

\mdfdefinestyle{mdgreenbox}{linecolor=ForestGreen,backgroundcolor=ForestGreen!5,
linewidth=2pt,rightline=false,leftline=true,topline=false,bottomline=false,}
\declaretheoremstyle[headfont=\bfseries\sffamily\color{ForestGreen!70!black},
mdframed={style=mdgreenbox},headpunct={ --- },]{thmgreenbox}

\mdfdefinestyle{mdorangebox}{linecolor=RedOrange,backgroundcolor=RedOrange!5,
linewidth=2pt,rightline=false,leftline=true,topline=false,bottomline=false,}
\declaretheoremstyle[headfont=\bfseries\sffamily\color{RedOrange!70!black},
mdframed={style=mdorangebox},headpunct={ --- },]{thmorangebox}

\mdfdefinestyle{mdblackbox}{linecolor=black,backgroundcolor=RedViolet!5!gray!5,
linewidth=3pt,rightline=false,leftline=true,topline=false,bottomline=false,}
\declaretheoremstyle[mdframed={style=mdblackbox}]{thmblackbox}

\declaretheoremstyle[spaceabove=6pt,spacebelow=6pt]{boldhead}
\declaretheoremstyle[spaceabove=6pt,spacebelow=6pt,headfont=\normalfont\itshape]{italichead}

\declaretheorem[style=boldhead,name=Problem]{problem}
\declaretheorem[style=boldhead,name=Problem,numbered=no]{problem*}
\declaretheorem[style=boldhead,name=Setup]{setup} % for config geo
\declaretheorem[style=boldhead,name=Example,sibling=problem]{example}
\declaretheorem[style=boldhead,name=$\bigstar$ Problem,sibling=problem]{niceprob}
\declaretheorem[style=boldhead,name=Theorem,sibling=problem]{theorem}
\declaretheorem[style=boldhead,name=Corollary,sibling=theorem]{corollary}
\declaretheorem[style=boldhead,name=Proposition,sibling=theorem]{prop}
\declaretheorem[style=boldhead,name=Lemma,numberwithin=problem]{lemma}
\declaretheorem[style=boldhead,name=Theorem,numbered=no]{theorem*}
\declaretheorem[style=boldhead,name=Lemma,numbered=no]{lemma*}
\declaretheorem[style=boldhead,name=Claim,numberwithin=problem]{claim}
\declaretheorem[style=boldhead,name=Claim,numbered=no]{claim*}
\declaretheorem[style=boldhead,name=Remark,numberwithin=problem]{remark}
\declaretheorem[style=boldhead,name=Remark,numbered=no]{remark*}
\declaretheorem[style=boldhead,name=Definition,sibling=theorem]{definition}
\declaretheorem[style=boldhead,name=Definition,numbered=no]{definition*}

\providecommand{\ol}{\overline}
\providecommand{\ceil}[1]{\left\lceil #1 \right\rceil}
\providecommand{\floor}[1]{\left\lfloor #1 \right\rfloor}
\newcommand{\dd}{\,\mathrm{d}}
\providecommand{\ul}{\underline}
\providecommand{\eps}{\varepsilon}
\providecommand{\half}{\frac{1}{2}}
\providecommand{\CC}{\mathbb C}
\providecommand{\FF}{\mathbb F}
\providecommand{\II}{\mathbb I}
\providecommand{\NN}{\mathbb N}
\providecommand{\QQ}{\mathbb Q}
\providecommand{\RR}{\mathbb R}
\providecommand{\ZZ}{\mathbb Z}
\providecommand{\ind}{\mathbbm 1}
\providecommand{\dg}{^\circ}
\providecommand{\ii}{\item}
\providecommand{\setdiff}{\smallsetminus}
\providecommand{\alert}[1]{{\sffamily\textbf{\textcolor{blue}{#1}}}}
\providecommand{\inv}{^{-1}}
\providecommand{\mb}[1]{\mathbf{#1}}
\newcommand{\what}{\widehat}

\newenvironment{soln}{\begin{proof}[Solution]}{\end{proof}}

\providecommand{\pow}{\operatorname{Pow}}
\providecommand{\vocab}[1]{{\textbf{\textcolor{ForestGreen}{#1}}}}
\providecommand{\lcm}{\operatorname{lcm}}
\providecommand{\abs}[1]{\left\lvert #1\right\rvert}
\providecommand{\smabs}[1]{\lvert #1\rvert}
\providecommand{\innerprod}[1]{\langle\!\langle #1\rangle\!\rangle}
\providecommand{\norm}[1]{\left\| #1\right\|}
\providecommand{\smnorm}[1]{\| #1\|}
\providecommand{\gr}{\operatorname{gr}}

\usepackage{mathrsfs}

% place title at top right
\setlength{\droptitle}{-8ex}
\pretitle{\begin{flushright}\Large\bfseries}
\posttitle{\par\end{flushright}}
\preauthor{\begin{flushright}}
\postauthor{\end{flushright}}
\predate{\begin{flushright}}
\postdate{\end{flushright}}

\title{MAT 549 HW01}
\author{\large Aditya Pahuja}
\date{\large Due 02/10}
\begin{document}
\maketitle

\begin{problem}
    Let $M$ be a nonempty positive-dimensional smooth manifold
    with or without boundary. Show that $\mathfrak X(M)$
    is infinite-dimensional.
\end{problem}

\begin{soln}
    By working in local coordinates, we may assume that $M = \RR^n$.
    For $k\in\NN$, let $f_k\colon M\to[0,\infty)$ such that
    $f\inv(0) = \RR^n - \ol B_{1/k}(0)$ (e.g. $f(x) = 1/k^2 - \abs{x}^2$
    for $x\in B_{1/k}(0)$ and $f(x) = 0$ otherwise);
    that is, the supports of the $f_k$ are a strictly
    decreasing sequence of sets. Define the vector fields
    $X_k = f_k\frac{\partial}{\partial x_1}$.

    Given the vector fields $X_{i_1}$, \dots, $X_{i_m}$
    with $i_1 < \cdots < i_m$ that are linearly dependent,
    consider a linear combination $a_1 X_{i_1} + \cdots + a_m X_{i_m}$ for
    nonzero real $a_1$, \dots, $a_m$.
    By construction, $X_{i_1}$ is nonzero at some points
    outside $\ol B_{1/i_2}(0)$ while all the other vector fields
    are zero outside this closed ball, so the linear combination
    is nonzero on $B_{1/i_1}(0) - \ol B_{1/i_2}(0)$.
    Thus any finite linear combination of the $X_k$ is nonzero,
    meaning $\{X_1,X_2,\dots\}$ is an infinite linearly independent
    subset of $\mathfrak X(M)$, so the space of vector fields on $M$
    is infinite-dimensional.
\end{soln}

\begin{problem}
    Let $M$ be a smooth manifold with boundary.
    Show that there exists a global smooth vector field on $M$
    whose restriction to $\partial M$ is everywhere inward-pointing,
    and one whose restriction to $\partial M$
    is everywhere outward-pointing.
\end{problem}

\begin{soln}
    Let $(U_\alpha,\varphi_\alpha)_{\alpha\in I}$
    be a collection of charts on $M$ whose domains cover $M$,
    and let $(\psi_\alpha)_{\alpha\in I}$ be a partition of unity
    subordinate to the covering $(U_\alpha)_{\alpha\in I}$.
    Then, for each $\alpha\in I$, define a vector field $X_\alpha$
    on $U_\alpha$ as follows: if $U_\alpha$ is contained
    in the interior of $M$, then set $X_\alpha\equiv 0$;
    otherwise, define $X_\alpha = \frac\partial{\partial x_n}$
    in local coordinates, so that $X_\alpha|_{\partial M\cap U_\alpha}$
    is an inward-pointing vector at each point of
    $\partial M\cap U_\alpha$. Then,
    \begin{equation*}
	X = \sum_{\alpha\in I}\varphi_\alpha X_\alpha
    \end{equation*}
    is a global vector field for which $X|_{\partial M}(p)$
    is the sum of some inward-pointing vectors, hence inward-pointing,
    for every $p\in\partial M$.
    Negating $X$ yields a global vector field
    that is outward-pointing at every point of $\partial M$.
\end{soln}

\begin{problem}
    For each of the following vector fields on the plane,
    compute its coordinate representation in
    polar coordinates on the right half-plane.
    \begin{enumerate}[(a)]
	\ii $X = x\frac{\partial}{\partial x}
	+ y\frac{\partial}{\partial y}$
	\ii $Y = x\frac{\partial}{\partial y}
	- y\frac{\partial}{\partial x}$
	\ii $Z = (x^2 + y^2)\frac{\partial}{\partial x}$.
    \end{enumerate}
\end{problem}

\begin{soln}
    By the change of coordinates formula (equation (3.11)),
    \begin{equation*}
	\frac\partial{\partial r}
	= \cos\theta\cdot\frac\partial{\partial x}
	+ \sin\theta\cdot\frac\partial{\partial y},\qquad
	\frac\partial{\partial \theta}
	= -r\sin\theta\cdot\frac\partial{\partial x}
	+ r\cos\theta\cdot\frac\partial{\partial y}.
    \end{equation*}
    This implies that
    \begin{equation*}
        X = r\cos\theta\cdot\frac\partial{\partial x}
	+ r\sin\theta\cdot\frac\partial{\partial y}
	=\boxed{ r\frac{\partial}{\partial r}},\qquad
	Y = -r\sin\theta\cdot\frac\partial{\partial x}
	+ r\cos\theta\cdot\frac\partial{\partial y}
	=\boxed{ \frac\partial{\partial\theta}}
    \end{equation*}
    and we can solve the system of equations
    for $\frac\partial{\partial x}$ directly to get
    \begin{equation*}
	Z = \boxed{r^2\cos\theta\cdot\frac\partial{\partial r}
	- r\sin\theta\cdot\frac\partial{\partial\theta}}.\qedhere
    \end{equation*}
\end{soln}

\begin{problem}
    Show that there is a smooth vector field on $S^2$
    that vanishes at exactly one point.
\end{problem}

\begin{soln}
    Consider the vector field $X\in\mathfrak X(S^2)$ given by
    \begin{equation*}
	X(x,y,z) = \frac 12(y^2-x^2 + (1-z)^2, -2xy, 2x(1-z)).
    \end{equation*}
    The component functions are smooth on $\RR^3$,
    hence of course smooth on $S^2$.
    The vectors are tangent to $S^2$ because
    \begin{align*}
        X(x,y,z)\cdot(x,y,z)
	&= \half(xy^2 - x^3 + x(1-z)^2 - 2xy^2 + 2xz(1-z))\\
	&= \half(-x^3 - xy^2 + x(1-z^2))\\
	&= \frac x2(1-x^2 - y^2 - z^2) = 0.
    \end{align*}
    Finally, to see that the vector field is only zero at one point,
    observe that $X(x,y,z) = 0$ implies that $x = 0$ or $y = 0$
    by looking at the second component; the first case means
    the first component $y^2 + (1-z)^2 = 0$,
    hence $y=0$ and $1-z=0$; on the other hand,
    the second case (assuming $x\ne 0$)
    implies that $1-z=0$ from the third component
    and then $x = 0$ because the first component equals $x^2 = 0$.
    In either case, $X(x,y,z) = 0$ implies $(x,y,z) = (0,0,1)$,
    so $X$ is nonzero away from $(x,y,z) = (0,0,1)$.
\end{soln}

\begin{problem}
    For each of the following pairs of vector fields,
    compute the Lie bracket $[X,Y]$.
    \begin{enumerate}[(a)]
	\ii $X = y\frac{\partial}{\partial z}
	- 2xy^2\frac{\partial}{\partial y}$;
	$Y = \frac{\partial}{\partial y}$
	\ii $X = x\frac\partial{\partial y} - y\frac\partial{\partial x}$;
	$Y = y\frac\partial{\partial z}-z\frac\partial{\partial y}$
	\ii $X = x\frac\partial{\partial y} - y\frac\partial{\partial x}$;
	$Y = x\frac\partial{\partial y} + y\frac\partial{\partial x}$
    \end{enumerate}
\end{problem}

\begin{soln}
    (a) \begin{equation*}
	[X,Y] = X(1)\frac\partial{\partial y}
	- Y(-2xy^2)\frac\partial{\partial y}
	- Y(y)\frac\partial{\partial z}
	= \boxed{4xy\frac\partial{\partial y} - \frac\partial{\partial z}}
    \end{equation*}

    (b) \begin{equation*}
	[X,Y] = X(-z)\frac\partial{\partial y}
	+ X(y)\frac\partial{\partial z}
	- Y(-y)\frac\partial{\partial x}
	- Y(x)\frac\partial{\partial y}
	= \boxed{x\frac\partial{\partial z} - z\frac\partial{\partial x}}
    \end{equation*}

    (c) \begin{equation*}
	[X,Y] = X(x)\frac\partial{\partial y}
	+ X(y)\frac\partial{\partial x}
	- Y(-y)\frac\partial{\partial x}
	- Y(x)\frac\partial{\partial y}
	= \boxed{2x\frac\partial{\partial x} - 2y\frac\partial{\partial y}}
	\qedhere
    \end{equation*}
\end{soln}











\end{document}
